\subsection{The radical-ideal frame}
\label{sec:annex-p10-radical-ideal-spectrum}

\begin{annexitem}{annex:p10-radical-spectrum-frame}{thm:p10-radical-spectrum-frame}{Frame laws}
\begin{proof}
An intersection of ideals is an ideal. If the ideals in the intersection
are radical and \(a^2\) belongs to all of them, then \(a\) belongs to
all of them, so the intersection is radical.
Consequently \(\rho(U)\) is the least radical ideal containing \(U\).
It follows that intersections and radical closures of unions have
the asserted meet and join universal properties. The empty intersection
is \(A\), and the empty join is \(\rho(\varnothing)\).

We first establish a useful closure calculation. If \(K\) is a
radical ideal and \(R\) is any subobject of \(A\), define
\[
 (K:R)=\{a\in A:\ \forall r\in R,\ ar\in K\}.
\]
It is an ideal: zero satisfies the condition, addition and additive
inverses follow from those of \(K\), and multiplication by a ring
element preserves it. It is radical as well. If \(a^2\in(K:R)\),
then \(a^2r\in K\) for every \(r\in R\). Hence
\((ar)^2=(a^2r)r\in K\), so \(ar\in K\), since \(K\) is radical.
Thus \(a\in(K:R)\).

To prove frame distributivity, let \(R\) be a radical ideal,
let \(I_j\) be any internally indexed family of radical ideals,
and put \(K=\bigvee_j(R\cap I_j)\).
For \(s\in I_j\) and \(r\in R\), the product \(rs\) belongs
to both ideals, hence to \(K\).
Therefore \(I_j\subseteq(K:R)\) for every \(j\).
The latter ideal is radical, so \(\bigvee_j I_j\subseteq(K:R)\).
If \(a\in R\cap\bigvee_j I_j\), this inclusion gives \(a^2\in K\),
and radicality gives \(a\in K\).
The opposite inclusion follows from monotonicity of meet and join.
We have proved
\[
 R\cap\bigvee_j I_j=\bigvee_j(R\cap I_j),
\]
including the empty-family case. Thus the complete lattice is a frame.

Every ideal contains zero, so \(D(0)=\rho(\varnothing)\).
An ideal containing one is all of \(A\), so \(D(1)=A\).
For the product identity, put \(K=D(ab)\).
Since \(ab\in K\), the radical ideal \((K:\{b\})\) contains
\(a\), hence contains \(D(a)\).
For every \(x\in D(a)\), it follows that \(xb\in K\).
The radical ideal \((K:\{x\})\) then contains \(b\), hence
contains \(D(b)\).
Thus \(xy\in K\) whenever \(x\in D(a)\) and \(y\in D(b)\).
In particular, if \(z\) belongs to both principal radical ideals,
then \(z^2\in K\), so \(z\in K\).
This proves \(D(a)\cap D(b)\subseteq D(ab)\).
Conversely, \(D(a)\) and \(D(b)\) are ideals containing \(ab\),
so \(D(ab)\) is contained in each.
The join \(D(a)\vee D(b)\) is an ideal containing \(a+b\);
leastness gives the asserted addition inequality.

Finally \(D(a)\subseteq I\) whenever \(a\in I\) and \(I\)
is radical. Hence \(\bigvee_{a\in I}D(a)\subseteq I\).
For the reverse inclusion, each \(a\in I\) belongs to its own
\(D(a)\), so belongs to that join. This proves generation.
Every step is valid in intuitionistic internal logic.
\end{proof}
\end{annexitem}

\begin{annexitem}{annex:p10-radical-spectrum-presentation}{thm:p10-radical-spectrum-presentation}{The universal frame map}
\begin{proof}
For \(v\in H\), let \(J_v=\{a\in A:d(a)\leq v\}\).
This is a radical ideal. It contains zero because \(d(0)=0\).
Addition closure follows from
\(d(a+b)\leq d(a)\vee d(b)\).
Multiplication by arbitrary \(r\) preserves membership since
\(d(ra)=d(r)\wedge d(a)\leq d(a)\).
Also
\(d(-1)=d(-1)\wedge d(-1)=d(1)=1\), so
\(d(-a)=d(-1)\wedge d(a)=d(a)\); additive inverses are preserved.
Finally \(d(a^2)=d(a)\), which proves radicality.

Define \(\varphi(I)=\bigvee_{a\in I}d(a)\).
This assignment is monotone.
It sends \(D(a)\) to \(d(a)\):
one inequality uses \(a\in D(a)\); the other follows because
\(J_{d(a)}\) is radical and contains \(a\), so contains \(D(a)\).

For an arbitrary family \(I_j\), put \(v=\bigvee_j\varphi(I_j)\).
Every \(I_j\) is contained in the radical ideal \(J_v\).
Hence their join is contained in \(J_v\), and
\(\varphi(\bigvee_j I_j)\leq v\).
The converse inequality follows by monotonicity.
Thus \(\varphi\) preserves arbitrary joins, including the empty join.

It preserves the top since \(1\in A\) and \(d(1)=1\).
For binary meets, frame distributivity in \(H\) gives
\[
 \begin{aligned}
 \varphi(I)\wedge\varphi(J)
 &=\bigvee_{\substack{a\in I\\b\in J}}
       \bigl(d(a)\wedge d(b)\bigr)\\
 &=\bigvee_{\substack{a\in I\\b\in J}}d(ab)
 \leq\varphi(I\cap J),
 \end{aligned}
\]
because \(ab\in I\cap J\).
The reverse inequality follows from the two inclusions of the
intersection. Hence \(\varphi\) is a frame homomorphism.

Any frame homomorphism with the specified generator values must
send \(I=\bigvee_{a\in I}D(a)\) to
\(\bigvee_{a\in I}d(a)\). This proves uniqueness.
When \(H\) is the truth-value frame, an arrow \(d:A\to H\)
is the classifier of a subobject of \(A\).
Its four displayed relations are exactly the four prime-predicate
conditions. This proves the stated point correspondence without
assuming an existence theorem for such points.
\end{proof}
\end{annexitem}
